\section{Genuine DOMAIN2 StepOne representative values and entire cuts without canonicality}

The ambient dimension is positive. Relative equivalence means the actual
same carrier, ordered boundary and mark with genuinely equal relative
original germ ideals. Every true history position is retained.

\begin{theorem}[Only actual original relative-ideal geometry constructs all DOMAIN2 StepOne representative values]
\label{hir:original-bound-two-stepone-representative-values}
\uses{hir:original-relative-ideal-level-cut-transport,hir:original-relative-ideal-positive-head-values,hir:original-zero-contact-representative-values,hir:original-finite-level-cut-current-lower-tests,hir:step-one-representative}
\lean{Hironaka.HironakaOriginalBoundTwoStepOneRepresentative.step_one_values_of_restrictEquiv}
\leanok
Let an original marked ideal be in DOMAIN2 with actual maximal order, and
let another coherent original marked ideal be genuinely relative-equivalent
to it. On ANY actual original level-cut history, at EVERY point of its
actual current cosupport, their ENTIRE unchanged original StepOne values
agree at EVERY positive lower FULL depth k+1. The second DOMAIN2, maximal
order, whole level-cut recipe, lower TESTs and queried points are supplied
internally. No lower CAN, source history/cut or desired numerical equality
is a caller input.
\end{theorem}
\begin{proof}
Actual relative-ideal domain transport supplies DOMAIN2 for the coherent
second reference. Maximal order also follows internally: keep the genuine
relative ideals and replace both equal marks by their successors. Their
true controlled cosupports at the reflexive history are equal. The original
maximal-order condition says exactly that the first such successor-mark
cosupport is empty; hence the second is empty too. This proves genuine
second maximal order without supplying it as a desired hypothesis.

The proved level-cut transport theorem constructs the entire second
level-cut history using actual admissibility, cosupport, level and special
identities. The real current point transfers by the same whole cosupport
identity. Apply genuine finite-level-cut geometry independently to both
references to obtain their actually computed lower DOMAIN1, restricted
TESTs and mask points before any value comparison. Original relative-ideal
calculus identifies the two levels, special predicates and regimes on
every actual preceding history. Actual regime-run induction therefore
gives identical run lengths and the SAME genuine split head and tail.

At a special point substitute the true level and special identities in
the unchanged original special StepOne branch. At a nonspecial point use
the identical actual run split and the independently constructed ordinary
lower geometry. At positive level the proved whole positive-mask/member
relative equality gives the all-depth lower-value identity. At zero use
the proved local tangent transfer, actual nested curve geometry, whole
kept-contact relative ideals and independently chosen tangent gluing;
its ordinary lower DOMAIN1/TEST/mask hypotheses are supplied by the prior
current-lower constructions. Substitute these lower identities and the
true common level into the literal original nonspecial StepOne function.
No canonical segments or target comparison certificates occur.
\end{proof}

\begin{theorem}[Actual entire StepOne-cut histories transport to coherent relative representatives]
\label{hir:original-bound-two-stepone-representative-cut}
\uses{hir:original-bound-two-stepone-representative-values,hir:original-finite-stepone-cut-history,hir:marked-equivalence-basic}
\lean{Hironaka.HironakaOriginalBoundTwoStepOneRepresentative.step_one_cut_of_restrictEquiv}
\leanok
For the same original inputs, an entire actual original StepOne-cut history
at lower depth k+1 constructs an entire actual StepOne-cut history of the
second reference on the SAME actual tower. No second cut, canonicality,
centre-equals-maximum condition or lower geometry is supplied.
\end{theorem}
\begin{proof}
Induct on the real finite tower. The reflexive history is a cut history.
At a next constructor the original whole relative-ideal theorem separately
transfers true last-centre admissibility and entire preceding cosupport.
The induction hypothesis supplies the second preceding cut. The original
preceding cut supplies its real level-cut recipe by the unchanged first
coordinate. At each second cosupport point use the genuine cosupport
identity to obtain the original point; every actual centre point is also
in that old cosupport by actual original admissibility. The preceding
entire-value theorem therefore applies at BOTH points. Rewrite the old
whole-centre inequality to obtain the second one and combine it with
actual second admissibility and the induction hypothesis. This constructs
the actual next second cut. No numerical inequality is used as a substitute
for geometry, no centre is asserted equal to the entire maximum, and no
next CAN, noncanonical FULL companion-cut construction, higher-carrier
representative invariance or global driver termination is concluded.
\end{proof}
